% Part: first-order-logic
% Chapter: completeness
% Section: compactness

\documentclass[../../../include/open-logic-section]{subfiles}

\begin{document}

\iftag{FOL}
      {\olfileid{fol}{com}{com}}
      {\olfileid{pl}{com}{com}}

\olsection{संहतता प्रमेय}

संपूर्णता प्रमेयाची एक महत्त्वाची निष्पत्ती म्हणजे संहतता प्रमेय.
!!{sentence}s च्या एखाद्या संचाचा प्रत्येक \emph{सांत} उपसंच पूर्ततायोग्य असेल,
तर संपूर्ण संच पूर्ततायोग्य असतो---तो संच स्वतः अनंत असला तरीही, असे
संहतता प्रमेय सांगते. हे मुळीच उघड नाही. प्रथमदर्शनी तरी अशी शक्यता
नाकारली जाईल असे काही दिसत नाही की !!{sentence}s चे अनंत संच व्याघाती
असतील, पण व्याघात, जणू काही, अनंत संख्येमुळेच उद्भवत असेल. अशी परिस्थिती
नाकारता येते, असे संहतता प्रमेय सांगते: ज्यांचा प्रत्येक सांत उपसंच
पूर्ततायोग्य आहे असे !!{sentence}s चे अपूर्ततायोग्य अनंत संच नसतात.
संपूर्णता प्रमेयाप्रमाणे याची तार्किक निष्पन्नतेशी संबंधित आवृत्तीही आहे:
!!{sentence}s च्या अनंत संचापासून काही निष्पन्न होत असेल, तर ते आधीच
एखाद्या सांत उपसंचापासून निष्पन्न होते.

\begin{defn}
  !!{formula}s चा संच~$\Gamma$ \emph{सांततः पूर्ततायोग्य} असतो तेव्हा आणि
  केवळ तेव्हाच, जेव्हा प्रत्येक सांत $\Gamma_0 \subseteq \Gamma$
  पूर्ततायोग्य असतो.
\end{defn}

\begin{thm}[संहतता प्रमेय]
\ollabel{thm:compactness}
वाक्यांचा कोणताही संच~$\Gamma$ आणि कोणतेही वाक्य~$!A$ यांच्यासाठी पुढील
विधाने लागू होतात:
\begin{enumerate}
  \item $\Gamma \Entails !A$ तेव्हा आणि केवळ तेव्हाच, जेव्हा असा सांत
    $\Gamma_0 \subseteq \Gamma$ असतो की $\Gamma_0 \Entails !A$.
  \item $\Gamma$ पूर्ततायोग्य असतो तेव्हा आणि केवळ तेव्हाच, जेव्हा तो
    सांततः पूर्ततायोग्य असतो.
\end{enumerate}
\end{thm}

\begin{proof}
आपण (2) सिद्ध करू. $\Gamma$ पूर्ततायोग्य असेल, तर अशी
\iftag{FOL}{!!a{structure}~$\Struct{M}$}{!!a{valuation}~$\pAssign{v}$}
असते की प्रत्येक $!A \in \Gamma$ साठी
$\iftag{FOL}{\Sat{M}}{\pSat{v}}{!A}$. अर्थात, ही
$\iftag{FOL}{\Struct{M}}{\pAssign{v}}$ देखील~$\Gamma$ चा प्रत्येक सांत
उपसंच पूर्त करते; म्हणून~$\Gamma$ सांततः पूर्ततायोग्य आहे.

आता $\Gamma$ सांततः पूर्ततायोग्य आहे असे समजा. मग प्रत्येक सांत
उपसंच~$\Gamma_0 \subseteq \Gamma$ पूर्ततायोग्य आहे. निर्दोषतेनुसार
(\iftag{FOL}{%
  \tagrefs{prfAX/{fol:axd:sou:cor:consistency-soundness},
    prfSC/{fol:seq:sou:cor:consistency-soundness},
    prfND/{fol:ntd:sou:cor:consistency-soundness},
    prfTab/{fol:tab:sou:cor:consistency-soundness}}}{%
  \tagrefs{prfAX/{pl:axd:sou:cor:consistency-soundness},
    prfSC/{pl:seq:sou:cor:consistency-soundness},
    prfND/{pl:ntd:sou:cor:consistency-soundness},
    prfTab/{pl:tab:sou:cor:consistency-soundness}}}),
प्रत्येक सांत उपसंच सुसंगत आहे. मग पुढील संदर्भांनुसार $\Gamma$ स्वतः
सुसंगत असला पाहिजे:
\iftag{FOL}{%
  \tagrefs{prfAX/{fol:axd:ptn:prop:proves-compact},
    prfSC/{fol:seq:ptn:prop:proves-compact},
    prfND/{fol:ntd:ptn:prop:proves-compact},
    prfTab/{fol:tab:ptn:prop:proves-compact}}}{%
  \tagrefs{prfAX/{pl:axd:ptn:prop:proves-compact},
    prfSC/{pl:seq:ptn:prop:proves-compact},
    prfND/{pl:ntd:ptn:prop:proves-compact},
    prfTab/{pl:tab:ptn:prop:proves-compact}}}.
संपूर्णतेनुसार (\olref[cth]{thm:completeness}), $\Gamma$ सुसंगत
असल्यामुळे तो पूर्ततायोग्य आहे.
\end{proof}

\tagprob{FOL}
\begin{prob}
\olref[fol][com][com]{thm:compactness} मधील (1) सिद्ध करा.
\end{prob}
\tagendprob

\tagprob{notFOL}
\begin{prob}
\olref[pl][com][com]{thm:compactness} मधील (1) सिद्ध करा.
\end{prob}
\tagendprob

\iftag{FOL}{%
\begin{ex}
उपपत्ती~$\Gamma$ च्या प्रत्येक प्रतिमान~$\Struct{M}$ मध्ये प्रत्येक
पद~$t$ हे अर्थातच~$\Domain{M}$ मधील !!a{element} निर्देशित करते.
$\Domain{M}$ मधील प्रत्येक !!{element} देखील कोणत्यातरी पदाने निर्देशित
केला जातो, याची हमी देता येते का? दुसऱ्या शब्दांत, ज्यांची सर्व प्रतिमाने
बंद पदांनी आच्छादित आहेत अशा उपपत्ती~$\Gamma$ असतात का? $\Gamma$ ला
अनंत प्रतिमाने असतील, तर असे नसते, हे संहतता प्रमेय दाखवते. ते पुढीलप्रमाणे
दिसते: $\Struct{M}$ हे~$\Gamma$ चे अनंत प्रतिमान असू द्या आणि $c$ हे
$\Gamma$ च्या भाषेत नसलेले !!a{constant} असू द्या. $\Gamma$ ची
भाषा~$\Lang{L}$ मधील प्रत्येक पद~$t$ साठी असलेल्या सर्व
$\eq/[c][t]$ या वाक्यांचा~$\Delta$ हा संच असू द्या, म्हणजे
  \[
  \Delta = \Setabs{\eq/[c][t]}{t \in \Trm[L]}.
  \]
$\Gamma \cup \Delta$ चा सांत उपसंच $\Gamma' \cup \Delta'$ असा लिहिता
येतो; येथे $\Gamma' \subseteq \Gamma$ आणि $\Delta' \subseteq \Delta$.
$\Delta'$ सांत असल्यामुळे त्यात फक्त सांत अनेक पदे असू शकतात.
त्यांपैकी कोणत्याही पदाने निर्देशित न केलेला $\Domain{M}$ मधील
!!a{element} $a \in \Domain{M}$ असू द्या आणि $\Struct{M'}$ ही
$\Struct{M}$ सारखीच, पण $\Assign{c}{M'} = a$ देखील असलेली !!{structure}
असू द्या. $\Delta'$ मध्ये येणाऱ्या प्रत्येक~$t$ साठी
$a \neq \Value{t}{M}$ असल्यामुळे $\Sat{M'}{\Delta'}$. तसेच
$\Sat{M}{\Gamma}$, $\Gamma' \subseteq \Gamma$, आणि~$\Gamma$ मध्ये $c$
येत नसल्यामुळे $\Sat{M'}{\Gamma'}$ देखील. म्हणून $\Gamma \cup \Delta$
च्या प्रत्येक सांत उपसंच $\Gamma' \cup \Delta'$ साठी
$\Sat{M'}{\Gamma' \cup \Delta'}$. त्यामुळे $\Gamma \cup \Delta$ चा
प्रत्येक सांत उपसंच पूर्ततायोग्य आहे. संहततेनुसार $\Gamma \cup \Delta$
स्वतः पूर्ततायोग्य आहे. म्हणून $\Sat{M}{\Gamma \cup \Delta}$ अशी
प्रतिमाने असतात. असे प्रत्येक $\Struct{M}$ हे~$\Gamma$ चे प्रतिमान आहे,
पण ते बंद पदांनी आच्छादित नाही; कारण~$\Lang{L}$ मधील सर्व पदे~$t$
यांच्यासाठी $\Value{c}{M} \neq \Value{t}{M}$.
\end{ex}

\begin{ex}
!!{predicate}~$<$, !!{constant}s~$\Obj{0}$, $\Obj{1}$, आणि
!!{function}s~$+$, $\times$, व~$-$ असलेल्या !!a{language}~$\Lang{L}$ चा
विचार करा. प्रांत~$\Rat$ आणि उघड निर्वचनार्थ असलेल्या
!!{structure}~$\Struct{Q}$ मध्ये सत्य असलेल्या या !!{language} मधील सर्व
!!{sentence}s चा संच~$\Gamma$ असू द्या. $\Gamma$ म्हणजे परिमेय
संख्यांविषयी सत्य असलेल्या~$\Lang{L}$ मधील सर्व !!{sentence}s चा संच.
अर्थात~$\Rat$ मध्ये (आणि~$\Real$ मध्येही) अशा संख्या~$r$ नसतात की त्या
$0$ पेक्षा मोठ्या, पण प्रत्येक $k \in \PosInt$ साठी $1/k$ पेक्षा लहान
असतात. अशी संख्या अस्तित्वात असती, तर ती \emph{अतिसूक्ष्म संख्या} असती:
शून्येतर, पण अनंतपटीने लहान. $\Gamma$ ची अतिसूक्ष्म संख्या असलेली प्रतिमाने
अस्तित्वात आहेत, हे दाखवण्यासाठी संहतता प्रमेय वापरता येते. आपल्या भाषेत
भागाकाराचे !!a{function} नाही (शून्याने भागाकार अपरिभाषित असतो आणि
!!{function}s चे निर्वचन सर्वत्र परिभाषित फलनांनी केले पाहिजे). तरीही
$r < 1/k$ हे व्यक्त करता येते; कारण हे तेव्हा आणि केवळ तेव्हाच लागू होते,
जेव्हा $r \cdot k < 1$. आता $c$ हे नवीन !!{constant} असू द्या आणि
$\Delta$ पुढीलप्रमाणे असू द्या:
\[
\{\Obj{0} < c\}
\cup \Setabs{ c \times \num k < \Obj{1} }{k \in \PosInt}
\]
(येथे $\num{k} = (\Obj{1} + (\Obj{1} + \dots + (\Obj{1} +
\Obj{1})\dots))$ मध्ये $k$ इतके~$\Obj{1}$ आहेत). $\Delta$ च्या कोणत्याही
सांत उपसंच~$\Delta_0$ साठी असा~$K$ असतो की~$\Delta_0$ मधील सर्व
!!{sentence}s $c \times \num{k} < \Obj{1}$ यांमध्ये $k < K$ असते.
$\Assign{c}{Q'} = 1/K$ असे ठेवून $\Struct{Q}$ चा विस्तार
$\Struct{Q'}$ पर्यंत केला, तर कोणत्याही सांत $\Gamma_0 \subseteq \Gamma$
साठी $\Sat{Q'}{\Gamma_0 \cup \Delta_0}$, आणि म्हणून $\Gamma \cup \Delta$
सांततः पूर्ततायोग्य आहे (सराव: हे सविस्तर सिद्ध करा). संहततेनुसार
$\Gamma \cup \Delta$ पूर्ततायोग्य आहे. $\Gamma \cup \Delta$ च्या
कोणत्याही प्रतिमान~$\Struct{S}$ मध्ये एक अतिसूक्ष्म संख्या असते, ती म्हणजे
$\Assign{c}{S}$.
\end{ex}
}{}

\tagprob{FOL}
\begin{prob}
अंकगणिताच्या प्रमाण प्रतिमान~$\Struct{N}$ मध्ये प्रत्येक सूत्र~$\num{n}
< x$ ची पूर्ति करणारा कोणताही !!{element}~$k \in \Domain{N}$ नाही
(येथे $\num{n}$ हे $n$ इतकी~$\prime$ चिन्हे असलेले
$\Obj{0}^{\prime\dots\prime}$ आहे). अंकगणिताच्या प्रमाण
प्रतिमान~$\Struct{N}$ मध्ये सत्य असलेली अंकगणिताच्या भाषेतील
!!{sentence}s ही प्रत्येक सूत्र~$\num{n} < x$ ची पूर्ति \emph{करणारा}
!!a{element} असलेल्या !!a{structure}~$\Struct{N'}$ मध्येही सत्य आहेत,
हे दाखवण्यासाठी संहतता प्रमेय वापरा.
\end{prob}
\tagendprob

\iftag{FOL}{
\begin{ex}
!!{identity} असलेले प्रथम-कोटी तर्कशास्त्र हे !!{domain} च्या आकाराला
एखादी किमान मर्यादा असली पाहिजे, असे व्यक्त करू शकते हे आपल्याला माहीत
आहे: वाक्य~$!A_{\ge n}$ (जे ``किमान $n$ भिन्न वस्तू आहेत'' असे सांगते)
फक्त अशा रचनांमध्ये सत्य असते ज्यांत $\Domain{M}$ मध्ये किमान~$n$ वस्तू
असतात. म्हणून
  \[
  \Delta = \Setabs{!A_{\ge n}}{n \ge 1}
  \]
असे घेतले, तर~$\Delta$ चे कोणतेही प्रतिमान अनंत असले पाहिजे. म्हणून
एखाद्या उपपत्तीला फक्त अनंत प्रतिमानेच आहेत याची हमी तिच्यात~$\Delta$
मिळवून देता येते: $\Gamma \cup \Delta$ ची प्रतिमाने म्हणजे~$\Gamma$ ची
सर्व आणि फक्त अनंत प्रतिमाने.

अशा रीतीने प्रथम-कोटी तर्कशास्त्र अनंतता व्यक्त करू शकते. पण ते सांतता
व्यक्त करू शकत नाही, हे संहतता प्रमेय दाखवते. कारण वाक्यांचा
एखादा संच~$\Lambda$ सर्व आणि फक्त सांत !!{structure}s मध्ये पूर्त होतो
असे समजा. मग $\Delta \cup \Lambda$ सांततः पूर्ततायोग्य आहे. का? सांत
$\Delta' \cup \Lambda' \subseteq \Delta \cup \Lambda$ असा संच असू
द्या; येथे $\Delta' \subseteq \Delta$ आणि $\Lambda' \subseteq \Lambda$.
$!A_{\ge n} \in \Delta'$ असेल अशा संख्यांपैकी सर्वांत मोठी संख्या~$n$
असू द्या. $\Lambda$ ची सर्व सांत रचनांमध्ये पूर्ति होत असल्यामुळे,
त्याला सांत अनेक, पण $\ge n$ इतके !!{element}s असलेले
प्रतिमान~$\Struct{M}$ असते. पण मग $\Sat{M}{\Delta' \cup \Lambda'}$.
संहततेनुसार $\Delta \cup \Lambda$ ला अनंत प्रतिमान असते; पण हे
$\Lambda$ ची पूर्ति फक्त सांत !!{structure}s मध्ये होते या गृहीतकाच्या
विरुद्ध आहे.
\end{ex}
}{}

\end{document}
